\section{Шапка сайта}
\label{url:header}

\component{Верхняя панель}{cmp:root:header}

\paragraph{Назначение}
Глобальная навигация по разделам и доступ к профилю. Видна на всех страницах, кроме экрана игровой сессии.

\screenshot{components/header.png}{Хедер с навигацией}

\paragraph{Что видно без входа}
Неавторизованному пользователю доступны только логотип, пункт ``Главная'', раздел ``Контакты'' и кнопки ``Наблюдение'', ``Вход'', ``Регистрация''. Остальные пункты меню появляются после входа.

\action{Переход по логотипу}{act:root:header:logo}
\uilink{Логотип ``RPG Master''}{/}{ui:root:logo}
\paragraph{Действие}
Открывает главную страницу (\texttt{/}), см. \hyperref[url:root]{главу ``Главная''}.

\action{Переход: Главная}{act:root:header:nav-home}
\uilink{Пункт меню ``Главная''}{/}{ui:root:nav-home}
\paragraph{Действие}
Открывает главную страницу (\texttt{/}), см. \hyperref[url:root]{главу ``Главная''}.

\action{Переход: Сценарии}{act:root:header:nav-scenarios}
\uilink{Пункт меню ``Сценарии''}{/scenarios}{ui:root:nav-scenarios}
\paragraph{Действие}
Открывает список сценариев (\texttt{/scenarios}), см. \hyperref[url:scenarios]{главу ``Сценарии''}. Требует входа.

\action{Переход: Паки}{act:root:header:nav-entity-packs}
\uilink{Пункт меню ``Паки''}{/entity-packs}{ui:root:nav-entity-packs}
\paragraph{Действие}
Открывает наборы сущностей, переиспользуемые между сценариями (\texttt{/entity-packs}). Требует входа.

\action{Переход: Кампании}{act:root:header:nav-campaigns}
\uilink{Пункт меню ``Кампании''}{/campaigns}{ui:root:nav-campaigns}
\paragraph{Действие}
Открывает список кампаний (\texttt{/campaigns}) --- серий связанных подходов, см. раздел \ref{sec:concepts:chain}. Требует входа.

\action{Переход: Запущенные}{act:root:header:nav-launched}
\uilink{Пункт меню ``Запущенные''}{/launched-scenarios}{ui:root:nav-launched}
\paragraph{Действие}
Открывает список запущенных миров (\texttt{/launched-scenarios}). Требует входа.

\action{Переход: Группы}{act:root:header:nav-access-groups}
\uilink{Пункт меню ``Группы''}{/access_groups}{ui:root:nav-access-groups}
\paragraph{Действие}
Открывает управление доступом: группы и пользователи (\texttt{/access\_groups}), см. \hyperref[url:access-groups]{главу ``Группы''}. Содержимое доступно только администратору.

\action{Переход: Аудио}{act:root:header:nav-audio}
\uilink{Пункт меню ``Аудио''}{/audio}{ui:root:nav-audio}
\paragraph{Действие}
Открывает библиотеку музыки и звуков (\texttt{/audio}). Требует входа.

\action{Переход: Контакты}{act:root:header:nav-contacts}
\uilink{Пункт меню ``Контакты''}{/contacts}{ui:root:nav-contacts}
\paragraph{Действие}
Открывает страницу связи с командой и скачивания этого руководства (\texttt{/contacts}), см.\ \hyperref[url:contacts]{главу ``Контакты''}. Доступна без входа.

\action{Переход: Профиль}{act:root:header:nav-me}
\uilink{Имя пользователя и аватар}{/me}{ui:root:nav-me}
\paragraph{Действие}
Открывает профиль (\texttt{/me}), см. \hyperref[url:me]{главу ``Профиль''}. Корона рядом с именем означает права администратора.

\action{Наблюдение}{act:root:header:observer}
\uilink{Кнопка ``Наблюдение''}{/observer}{ui:root:btn-observer}
\paragraph{Действие}
Открывает вход в комнату наблюдателя по коду (\texttt{/observer}), см.\ главу \ref{ch:observer}. Показывается только неавторизованным.

\action{Выйти}{act:root:header:logout}
\uibutton{Кнопка ``Выйти''}{ui:root:btn-logout}
\paragraph{Действие}
Завершает сессию и возвращает на страницу входа.

\action{Открыть мобильное меню}{act:root:header:mobile-menu}
\uibutton{Кнопка меню (``бургер'')}{ui:root:btn-mobile-menu}
\paragraph{Действие}
Открывает боковое меню навигации на узких экранах.
